\chapter{R\'ealisation et analyse critique}
\finPageTitre

\section{Introduction}
Ce troisi\`eme chapitre pr\'esente concr\`etement les r\'ealisations techniques d\'evelopp\'ees tout au long du stage. Il d\'ecrit tout d'abord l'environnement de d\'eveloppement et la conteneurisation des services, d\'etaille ensuite l'impl\'ementation des pipelines d'ingestion et du bus de messages, expose les r\'esultats des bancs d'essai exp\'erimentaux, et conclut par une discussion critique des performances observ\'ees.

\section{Environnement de d\'eveloppement et infrastructure logicielle}
\subsection{Configuration des environnements et outillage}
Pour garantir la reproductibilit\'e des d\'eploiements, l'int\'egralit\'e des composants a \'et\'e configur\'ee sous forme de conteneurs Docker orchestr\'es par Docker Compose en local :
\begin{itemize}
    \item \textbf{Gestionnaire de versions} : D\'ep\^ot Git structur\'e avec validation automatis\'ee par hooks pr\'e-commit ;
    \item \textbf{Bus de messages} : Cluster Apache Kafka (version 3.6) coupl\'e \`a KRaft pour s'affranchir de ZooKeeper ;
    \item \textbf{Base de donn\'ees chronologique} : Instance TimescaleDB adoss\'ee \`a PostgreSQL 16.
\end{itemize}

\subsection{Organisation modulaire du code source}
Le projet adopte une architecture multi-d\'ep\^ots unifi\'ee en mono-repo avec des r\'epertoires d\'edi\'es pour chaque microservice (\texttt{services/gateway}, \texttt{services/processor}, \texttt{services/storage}). Chaque service encapsule sa propre logique de test et sa configuration d'environnement.

\section{Impl\'ementation des microservices principaux}
\subsection{Module d'ingestion : Gateway-Service}
Le service frontal d'ingestion r\'eceptionne les requ\^etes HTTP POST \'emises par les passerelles p\'eriph\'eriques IoT. D\`es r\'eception d'un payload JSON, le contr\^oleur valide la signature cryptographique et publie l'\'ev\'enement sur le topic Kafka \texttt{telemetry.raw}. La Figure~\ref{fig:module-1} illustre le tableau de bord de supervision de cette passerelle.

\begin{figure}[H]
\centering
\IfFileExists{figures/ch3/interface_module1.png}{%
    \includegraphics[width=0.85\textwidth]{figures/ch3/interface_module1.png}%
}{%
    \fbox{\parbox[c][4cm][c]{0.85\textwidth}{\centering \textit{[Interface de monitoring du microservice d'ingestion Gateway-Service]}}}%
}
\caption{Tableau de bord de monitoring du service d'ingestion en temps r\'eel}
\label{fig:module-1}
\end{figure}

L'interface visualis\'ee \`a la Figure~\ref{fig:module-1} permet d'observer en direct le d\'ebit instantan\'e de requ\^etes re\c{c}ues et le taux de succ\`es des publications dans le courtier.

\subsection{Module de traitement et calcul des m\'etriques}
Le microservice \texttt{Ingestion-Processor} consomme en continu les messages du topic \texttt{telemetry.raw}. Il applique un fen\^etrage temporel pour d\'etecter les anomalies thermiques ou vibratoires. La Figure~\ref{fig:module-2} montre l'architecture interne de ce pipeline de traitement.

\begin{figure}[H]
\centering
\IfFileExists{figures/ch3/interface_module2.png}{%
    \includegraphics[width=0.85\textwidth]{figures/ch3/interface_module2.png}%
}{%
    \fbox{\parbox[c][4cm][c]{0.85\textwidth}{\centering \textit{[Diagramme de flux interne du microservice de traitement Ingestion-Processor]}}}%
}
\caption{Diagramme de traitement asynchrone des flux d'\'ev\'enements}
\label{fig:module-2}
\end{figure}

L'examen de la Figure~\ref{fig:module-2} met en lumi\`ere le passage par une file d'attente d\'edi\'ee aux \'ev\'enements anormaux (\textit{Dead Letter Queue}), assurant qu'aucune donn\'ee malform\'ee ne bloque le pipeline principal.

\section{Mod\'elisation math\'ematique et r\'esultats de performance}
\subsection{Formulation de la latence de traitement}
La latence globale de bout en bout subie par une mesure t\'el\'em\'etrique est mod\'elis\'ee par l'\'equation~\eqref{eq:latence} :
\begin{equation}
L_{\text{total}} = L_{\text{r\'eseau}} + \frac{S_{\text{payload}}}{D_{\text{nominal}}} + T_{\text{buffer}} + T_{\text{traitement}}
\label{eq:latence}
\end{equation}
o\`u $L_{\text{r\'eseau}}$ repr\'esente le temps de transit r\'eseau, $S_{\text{payload}}$ la taille du paquet en octets, $D_{\text{nominal}}$ la bande passante disponible, $T_{\text{buffer}}$ le temps d'attente dans la file d'attente Kafka, et $T_{\text{traitement}}$ le temps d'ex\'ecution de la logique m\'etier. L'adoption du courtage asynchrone permet de stabiliser $T_{\text{buffer}}$ m\^eme en cas de surtension temporaire.

\subsection{Mesures exp\'erimentales sous charge de trafic}
Un g\'en\'erateur de charge a simul\'e l'activit\'e simultan\'ee de 5\,000 capteurs \'emettant \`a une fr\'equence de 10 Hertz. Le Tableau~\ref{tab:resultats-metriques} compare les performances mesur\'ees entre l'architecture initiale monolithique et la nouvelle architecture microservices.

\begin{table}[H]
\centering
\begin{tabularx}{\textwidth}{X c c c c}
\toprule
\textbf{Indicateur d'\'Evaluation} & \textbf{Syst\`eme Initial} & \textbf{Microservices} & \textbf{\'Evolution} & \textbf{Impact} \\
\midrule
\textbf{D\'ebit maximal soutenu}       & 1\,200 req/s              & 14\,500 req/s          & $\times 12$          & Gain majeur \\
\textbf{Latence m\'ediane (p50)}      & 340\,ms                   & 28\,ms                 & $-91{,}7\,\%$        & Fluidit\'e temps r\'eel \\
\textbf{Latence au 99\textsuperscript{e} percentile (p99)} & 2\,150\,ms & 95\,ms & $-95{,}6\,\%$ & Suppression des pics \\
\textbf{Taux de perte de trames}      & $4{,}8\,\%$               & $0{,}0\,\%$            & $-100\,\%$           & Z\'ero perte de donn\'ee \\
\bottomrule
\end{tabularx}
\caption{Synth\`ese comparative des indicateurs de performance sous charge}
\label{tab:resultats-metriques}
\end{table}

Les mesures consign\'ees dans le Tableau~\ref{tab:resultats-metriques} confirment l'efficacit\'e du partitionnement Kafka et de l'\'ecriture group\'ee dans TimescaleDB, le d\'ebit admissible \'etant multipli\'e par douze tout en abaissant drastiquement les temps de r\'eponse.

\section{Discussion critique et \'evaluation des limites}
\subsection{Gains op\'erationnels constat\'es}
L'architecture d\'eploy\'ee r\'esout le point critique de d\'efaillance unique. Lors d'un test d'arr\^et inopin\'e du service de stockage, les messages ont \'et\'e conserv\'es sans aucune perte dans Kafka, puis consomm\'es normalement d\`es le red\'emarrage du conteneur.

\subsection{Limites actuelles du banc exp\'erimental}
Bien que les r\'esultats d\'emontrent une nette am\'elioration, deux limites subsistent :
\begin{itemize}
    \item La gestion du cluster Kafka mono-n\oe ud en environnement local n\'ecessite une \'evolution vers un cluster multi-zones g\'eographiques pour r\'epondre aux contraintes de r\'esilience d'un environnement de production mondial ;
    \item L'absence actuelle d'un module d'agr\'egation analytique automatis\'e sur les flux historiques fige les requ\^etes complexes \`a des p\'eriodes inf\'erieures \`a trente jours.
\end{itemize}

\sectionTransition{%
Ce troisi\`eme chapitre a d\'etaill\'e l'impl\'ementation des microservices, pr\'esent\'e la formulation th\'eorique de la latence et quantifi\'e les gains de performance obtenus. Le quatri\`eme chapitre conclut ce rapport en dressant le bilan global des r\'ealisations et en formulant les perspectives techniques d'\'evolution.
}
